\section*{Hoofdstuk \ref{ch:g2:moon-in-the-sky} --- De Maan aan de hemel}

\begin{solution}{exo:g2:moon-in-the-sky:1}
Veel kleiner dan de Aarde, en veel verder weg dan de wolken --- verder
dan welk vliegtuig ook vliegt. We weten dat mensen haar kunnen
bezoeken omdat astronauten dat gedaan hebben: de reis kostte hun drie
dagen per keer.
\end{solution}

\begin{solution}{exo:g2:moon-in-the-sky:2}
Kraters --- komvormige littekens van rotsen uit de ruimte die op de
\omterm{def:g2:moon-in-the-sky:moon}{Maan} zijn ingeslagen. Ze blijven bestaan omdat de \omterm{def:g2:moon-in-the-sky:moon}{Maan} geen
\omterm{def:g2:air-around-us:wind}{wind} en geen regen heeft om ze te slijten.
\end{solution}

\begin{solution}{exo:g2:moon-in-the-sky:3}
Een lichtvanger. Maanlicht is zonlicht dat weerkaatst op het grijze
gesteente van de \omterm{def:g2:moon-in-the-sky:moon}{Maan}, zoals een witte muur een kamer
verlicht zonder een lamp te zijn.
\end{solution}

\begin{solution}{exo:g2:moon-in-the-sky:4}
Altijd de helft die naar de Zon toe gekeerd is --- de zonkant is
helder, de achterkant heeft haar eigen nacht.
\end{solution}

\begin{solution}{exo:g2:moon-in-the-sky:5}
Ja --- op veel middagen hangt er een bleke \omterm{def:g2:moon-in-the-sky:moon}{Maan} in de blauwe \omterm{def:g2:air-around-us:air}{lucht}.
Het beste bewijs voor een twijfelende vriend: ga op zo'n middag naar
buiten en wijs haar aan.
\end{solution}

\begin{solution}{exo:g2:moon-in-the-sky:6}
\omterm{ex:g2:moon-in-the-sky:shapes}{Maansikkel}, half, bijna vol, \omterm{ex:g2:moon-in-the-sky:shapes}{volle maan} --- en dan weer terug via
bijna vol, half en \omterm{ex:g2:moon-in-the-sky:shapes}{maansikkel} naar de verstopte ``nieuwe'' \omterm{def:g2:moon-in-the-sky:moon}{Maan}. De
hele voorstelling duurt ongeveer een maand.
\end{solution}

\begin{solution}{exo:g2:moon-in-the-sky:7}
Een schreeuw zou geen geluid maken --- geluid heeft \omterm{def:g2:air-around-us:air}{lucht} nodig. Een
vlieger zou niet kunnen vliegen: geen \omterm{def:g2:air-around-us:air}{lucht} betekent geen \omterm{def:g2:air-around-us:wind}{wind}, en
niets waar de vlieger op kan leunen.
\end{solution}

\begin{solution}{exo:g2:moon-in-the-sky:8}
Elke avond hetzelfde uur en dezelfde eerlijke tekening (dezelfde plek
helpt ook): een eerlijk dagboek verandert maar één ding tegelijk ---
de nacht.
\end{solution}

\begin{solution}{exo:g2:moon-in-the-sky:9}
De huizen zijn dichtbij, dus door erlangs te rijden verandert hoe Rosa
ze ziet --- ze glijden voorbij en zijn weg. De \omterm{def:g2:moon-in-the-sky:moon}{Maan} is zo ver dat de
hele rit haar uitzicht erop helemaal niet verandert: ze blijft naast
haar schouder en lijkt mee te gaan. De kampioen ver weg ``volgt''
altijd.
\end{solution}

\begin{solution}{exo:g2:moon-in-the-sky:10}
De Aarde zelf. Zonlicht kaatst van de heldere Aarde op de donkere kant
van de \omterm{def:g2:moon-in-the-sky:moon}{Maan} en verlicht die zwakjes --- aardschijnsel: onze planeet
die de \omterm{def:g2:moon-in-the-sky:moon}{Maan} een wederdienst bewijst.
\end{solution}
